% TOP_PROBLEM: {"id": 3053, "title": "The Double Cap Conjecture", "category_id": 2, "category": {"id": 2, "name": "combinatorics", "display_name": "Combinatorics", "description": "Counting problems, graph theory, discrete structures.", "slug": "combinatorics", "order_index": 2, "created_at": "2026-07-31T15:26:25.671Z"}, "status": "open", "problem_number": "OPG-60000", "set_id": 12}
% FIRST_POSTED: 2026-10-01
% ATTEMPT_STATUS: unresolved
% UnsolvedMath ID 3053; source catalogue code OPG-60000
% !TeX program = lualatex
% Research attempt by GPT-6 Astra Ultra; no claim of a new complete solution.
\documentclass[11pt,a4paper]{article}
\usepackage[margin=25mm]{geometry}
\usepackage{fontspec}
\setmainfont{lmroman10-regular.otf}[BoldFont=lmroman10-bold.otf,
 ItalicFont=lmroman10-italic.otf,BoldItalicFont=lmroman10-bolditalic.otf]
\usepackage{amsmath,amssymb,mathtools}
\usepackage{unicode-math}
\setmathfont{latinmodern-math.otf}
\AtBeginDocument{\let\setminus\smallsetminus}
\usepackage{xurl,hyperref}
\hypersetup{colorlinks=true,urlcolor=blue}
\setcounter{secnumdepth}{0}
\setlength{\parindent}{0pt}
\setlength{\parskip}{5pt}
\setlength{\emergencystretch}{3em}
\begin{document}
\section*{The Double Cap Conjecture}\label{3053}
{\small UnsolvedMath ID 3053 · OPG-60000 · Combinatorics · OpenGarden}
\subsection{Definitions and mathematical statement}
For an integer $n\ge2$, let
$S^{n-1}=\{x\in\mathbb R^n:\|x\|_2=1\}$, equipped with normalized
rotation-invariant surface measure $\sigma_{n-1}$, so that
$\sigma_{n-1}(S^{n-1})=1$. A measurable set $A\subseteq S^{n-1}$ avoids
orthogonal pairs if $\langle x,y\rangle\ne0$ for every $x,y\in A$.
This condition is pointwise: exceptions of measure zero are not allowed.

For a unit vector $u$, the union of the two opposite open caps of angular
radius $\pi/4$ is
\[
 D_u=\{x\in S^{n-1}:|\langle x,u\rangle|>1/\sqrt2\}.
\]
Vectors in the same cap make an angle less than $\pi/2$, while vectors
in opposite caps make an angle greater than $\pi/2$, so $D_u$ avoids
orthogonal pairs. Its measure depends on $n$ but not on the choice of $u$.

The conjecture asserts that for every measurable orthogonality-avoiding
$A\subseteq S^{n-1}$,
\[
 \sigma_{n-1}(A)\le\sigma_{n-1}(D_u).
\]
Equivalently, the maximum possible measure is achieved by this double-cap
construction. The statement concerns surface measure in each fixed
dimension; it is not a count of points or Euclidean volume in the ambient
space. Although cap boundaries have measure zero, open caps are used because
including all boundary points introduces orthogonal pairs. The problem
asks for the sharp extremal value, and does not assert uniqueness of every
maximizing measurable set.
\subsection{Short English statement}
Do two opposite open caps of angular radius forty-five degrees have the largest possible spherical measure among sets with no orthogonal pair?
\subsection{Research attempt}
\paragraph{Scope and process.}
This is a GPT-6 Astra Ultra research attempt dated 1 October 2026, using
primary-source browsing and elementary mathematical checks. The canonical
definition above is retained. Results below are restricted results or
reductions, not a claim of a new solution to the entire catalogue question.
No formal proof assistant or external mathematical referee checked this text.


\paragraph{Literature and plan.}
DeCorte and Pikhurko [R1] study upper bounds for orthogonality-avoiding
sets on spheres and discuss the double-cap conjecture. The elementary
orthonormal-frame bound below is the classical benchmark, not a new
sharp theorem. We prove it, verify the double caps pointwise, and
identify exactly the dimension-two equality and the higher-dimensional
gap left by this approach.

\paragraph{Antipodal symmetrization is legitimate.}
If $A$ avoids orthogonal pairs, then $A\cup(-A)$ does too. Indeed,
any putative orthogonal pair in this union can be given signs so that
both vectors lie in $A$, and changing signs preserves a zero inner
product. The union is measurable and has measure at least that of
$A$. Thus an extremal upper bound may be sought among antipodally
symmetric sets, without discarding a more dangerous nonsymmetric
competitor. This operation is an inclusion enlargement, not a claim
that arbitrary spherical rearrangements preserve orthogonality avoidance.

\paragraph{A universal upper bound.}
Choose a uniformly random orthonormal basis $(u_1,\ldots,u_n)$ of
$\mathbb R^n$, for example the columns of a Haar-distributed orthogonal
matrix. Pointwise orthogonality avoidance implies
\[
 \sum_{j=1}^n\mathbf1_A(u_j)\le1
\]
for every frame. Each column is uniform on $S^{n-1}$ by rotational
invariance. Integrating the displayed inequality gives
\[
 n\sigma_{n-1}(A)\le1.
\]
Measurability suffices for the expectation; no continuity or positive
distance from the forbidden inner product is required. The pointwise
condition on $A$ is important in justifying the frame inequality.

\paragraph{The lower construction and its precise boundary.}
After possibly replacing $x,y$ by their negatives, two points in
$D_u$ have positive $u$-coordinates $s,t>1/\sqrt2$. Write
$x=su+x_\perp$ and $y=tu+y_\perp$. Cauchy--Schwarz gives
\[
 \langle x,y\rangle\ge st-\sqrt{1-s^2}\sqrt{1-t^2}>0,
\]
because $s^2+t^2>1$, equivalently
$s^2t^2>(1-s^2)(1-t^2)$. Reversing either sign makes the inner
product negative but still nonzero. Thus $D_u$ avoids orthogonal
pairs. At the boundary, the vectors $(u+v)/\sqrt2$ and
$(u-v)/\sqrt2$ for $v\perp u$ are orthogonal, so replacing the
strict cap inequality by a weak one would fail pointwise.

\paragraph{Sharpness in dimension two and the gap in dimension three.}
On $S^1$ the two open caps are two arcs each of angular length
$\pi/2$, so their normalized measure is $1/2$. The frame upper
bound is also $1/2$, proving the conjectured extremal value for
$n=2$. Including endpoints is unnecessary to attain this measure.

On $S^2$, the normalized area of a cap with $u$-coordinate greater
than $a$ is $(1-a)/2$. This follows directly from the surface element
$d\theta\,dz$ in height coordinates: total area is $4\pi$ and the
cap area is $2\pi(1-a)$. Hence the double-cap measure is
\[
 1-\frac1{\sqrt2}\approx0.292893,
\]
whereas the frame bound is $1/3\approx0.333333$. The latter exceeds
the conjectured value by $1/\sqrt2-2/3>0$, since squaring the
positive quantities gives $1/2>4/9$. Thus the frame argument is
provably insufficient already in the first higher-dimensional case.

\paragraph{Verification and failed extension.}
The fact that the two-dimensional argument is sharp cannot be
integrated over great circles to recover the desired cap measure
in dimension three. Such an averaging only bounds every circle
section by one half and gives the weaker global upper bound one
half. Orthogonality constraints between points on different circles
are lost. The random-frame method retains more of those constraints
and improves the bound to $1/n$, but still does not encode the
geometry needed for the double-cap value. The first remaining gap
is a sharper all-measurable-set upper bound for $n\ge3$, not a
correction to the explicit cap construction.

\paragraph{Final status.}
\textbf{UNRESOLVED} in arbitrary dimension. Antipodal reduction,
the frame bound, pointwise admissibility of the construction, and
the exact dimension-two optimum are proved; no general sharp bound
or uniqueness assertion is claimed.

\subsection{Sources}
\begin{itemize}
\item[S1] ULAM.AI and the UnsolvedMath Contributors, supplied problems.json, record 3053 (OPG-60000), OpenGarden collection. Catalogue identity and source transcription; CC BY 4.0.
\url{https://huggingface.co/datasets/ulamai/UnsolvedMath}.
\item[S2] Open Problem Garden, The Double Cap Conjecture. Original problem and discussion; the mathematical formulation below is expanded from the supplied source record.
\url{http://www.openproblemgarden.org/op/the_double_cap_conjecture}.

\item[R1] Evan DeCorte and Oleg Pikhurko, \emph{Spherical Sets
Avoiding a Prescribed Set of Angles}, International Mathematics
Research Notices, article rnv319 (2015). Introduction checked
for Witsenhausen's frame bound and the double-cap question.
\url{https://wrap.warwick.ac.uk/id/eprint/76069/2/WRAP_Int%20Math%20Res%20Notices-2015-DeCorte-imrn-rnv319.pdf}.

\end{itemize}
\end{document}
